% ─────────────────────────────────────────────────────────────────────────────
% The Harbor Volume Workbook — Part I: The Single-Writer Kernel (questions)
%
% Extracted from the paper "The Single-Writer Kernel" (Floor 1, Kernel).
% Renumbered locally 1..n. Three tiers — Check / Trace / Open — with a star
% glyph for the open research problems. Open problems carry their global OP tag
% so they cross-link the paper's consolidated Open-Problems list.
%
% This file expects the workbook preamble to define:
%   \wbq{<tier>}{<label>}{<body>}   one question, small-caps tier eyebrow
%   \wbstar                          the difficulty glyph (filled star)
% A self-contained fallback definition is provided below if undefined.
% ─────────────────────────────────────────────────────────────────────────────

\providecommand{\wbstar}{\ensuremath{\bigstar}}
\providecommand{\wbq}[3]{%
  \par\vspace{6pt}\noindent
  {\scshape\bfseries\textcolor{hhcobalt}{#1.}}\enspace{\footnotesize\textcolor{hhgray}{#2}}\par
  \nopagebreak\vspace{2pt}\noindent #3\par\vspace{2pt}%
}

\section{The Single-Writer Kernel}
\label{wb:swk}

\noindent\emph{Floor 1 --- Kernel. Source paper: \emph{The Single-Writer Kernel}.
The \textsc{check} tier confirms you read; \textsc{trace} asks you to walk a
mechanism step by step; \textsc{open} is live research, starred and cross-tagged to
the paper's Open-Problems list.}

% ─── 1–3 · The kernel as seven organs (paper §3) ──────────────────────────────
\subsection*{The seven organs}

\wbq{Check}{Q1}{Which of the seven organs would still be meaningful if the daemon
ran over an in-memory database with no disk at all? Which become vacuous?}

\wbq{Trace}{Q2}{Pick any one organ and name the single SQLite table that is its
``home'' table. Then name one invariant that organ owes the layer above it.}

\wbq{Open \wbstar}{Q3}{The seven-organ decomposition is a choice. Is there a more
natural cut --- five organs, or nine? Argue for an alternative and show what it
makes easier to reason about and what it hides.}

% ─── 4–6 · The substrate organ (paper §4) ─────────────────────────────────────
\subsection*{The substrate organ}

\wbq{Check}{Q4}{The schema-by-union design lets any module create its tables in any
order. Give a concrete two-module example where a lazy column-rename migration would
be unsafe under this design.}

\wbq{Trace}{Q5}{Follow a single claim request from command line to disk. At how many
points does the single-writer discipline (Definition: single-writer discipline) hold
by \emph{routing} versus by the \emph{file lock}?}

\wbq{Open \wbstar}{Q6}{\textbf{(OP-7, Schema evolution without a sequencer.)} Can
idempotent self-init be extended to safe renames, backfills, and down-migrations
while preserving ``any module, any order''? Or is a version table unavoidable?}

% ─── 7–9 · Durability, split by fault class (paper §5) ────────────────────────
\subsection*{Durability by fault class}

\wbq{Check}{Q7}{Classify each fault as I1a-survivable or I1b-exposed: (a) a forced
kill of the daemon process; (b) a clean operating-system reboot; (c) yanking the
power cord; (d) an unhandled exception in a request handler.}

\wbq{Trace}{Q8}{A claim commits at $t_0$; the next auto-checkpoint fires at
$t_0+\delta$. Draw the window during which a power loss reverts the claim, and state
how a page-count auto-checkpoint threshold bounds $\delta$.}

\wbq{Open \wbstar}{Q9}{\textbf{(OP-10, Per-write-path durability.)} Design a
per-write-path durability selector: how would the kernel let a settlement-ledger
write opt into I1b while keeping claim writes fast? What is the cleanest interface
that does not re-introduce the durability/consistency conflation?}

% ─── 10–12 · The resource organ (paper §6) ────────────────────────────────────
\subsection*{The resource organ}

\wbq{Check}{Q10}{Why does releasing a lock you do not hold return success rather
than an error? Which property (I2--I4) is this, and what would break if it raised?}

\wbq{Trace}{Q11}{Two agents call \texttt{acquire} on the same fresh key at the same
instant. Walk the uniqueness-constraint insert for both and show exactly where the
loser learns the holder's identity.}

\wbq{Open \wbstar}{Q12}{\textbf{(OP-1, Fair exclusion without a scheduler.)} Can the
substrate add a bounded-wait fairness guarantee to claims and locks while keeping the
single-writer, no-background-scheduler simplicity? Is a FIFO wait-list of
time-to-live'd reservations enough, or does fairness force a scheduler --- and thus
blur the substrate/coordination boundary?}

% ─── 13–15 · The communication organ (paper §7) ───────────────────────────────
\subsection*{The communication organ}

\wbq{Check}{Q13}{Why does read-time marker decay make the system more honest than a
background tick alone? Give a scenario where tick-only decay would report a stale
signal as fresh.}

\wbq{Trace}{Q14}{A request speech act is sent over the bus and redelivered once by
the at-least-once channel. Trace what the operator should see under the silent-dedup
resolution, and what an envelope idempotency key would change.}

\wbq{Open \wbstar}{Q15}{\textbf{(OP-8, Stigmergic decay calibration.)} What per-kind
marker decay rate keeps stigmergic signals neither rotten nor prematurely
evaporated? Frame this as an empirical measurement, not a constant to guess.}

% ─── 16–19 · The obligation & enforcement organ (paper §8) ────────────────────
\subsection*{The obligation \& enforcement organ}

\wbq{Check}{Q16}{For each monitor rule in the post-commit enforcement theorem, say
whether it could in principle be moved to a pre-commit check (a before-insert
trigger). Which genuinely cannot, and why? (Hint: heartbeat freshness is a failure
detector, and Chandra and Toueg show failure detection cannot be made perfectly
synchronous.)}

\wbq{Trace}{Q17}{An agent attempts to close a commitment with an empty oracle
reference. Walk the Law-2 gate and show exactly where and why the transition is
refused.}

\wbq{Open \wbstar}{Q18}{\textbf{(OP-2, Regimentation vs.\ enforcement boundary.)}
Formalize the boundary: which monitor rules are truly regimentable (rejectable at
insert) versus only enforceable (detectable after)? A clean theorem here is the
deontic heart of the layer.}

\wbq{Open \wbstar}{Q19}{\textbf{(OP-5, Oracle completeness.)} Extend the Law-2 oracle
vocabulary to capture a ``done'' condition it currently misses, without re-admitting
free text. State your new oracle and prove it is Goodhart-resistant.}

% ─── 20–22 · The continuity organ (paper §9) ──────────────────────────────────
\subsection*{The continuity organ}

\wbq{Check}{Q20}{Why is ``recovery passes notes'' fundamentally weaker than
``recovery restores work'' for a language-model agent that died mid-edit? Name one
thing notes preserve and one thing they cannot.}

\wbq{Trace}{Q21}{Follow the note-monotonicity rule from an agent appending a note to
the monitor detecting a regression. At which point is the violation already durable
(per the post-commit enforcement theorem)?}

\wbq{Open \wbstar}{Q22}{\textbf{(OP-4, Checkpoint with teeth.)} What is the minimal
durable artifact that turns ``recovery passes notes'' into ``recovery restores
work''? Is it a content-addressed snapshot of \{working-tree diff $+$ open claims $+$
commitment set $+$ the last $N$ turns of the agent's transcript\}? What is recoverable
versus fundamentally lost when a language-model agent dies mid-thought?}

% ─── 23–25 · The invariants as theorems (paper §11) ───────────────────────────
\subsection*{The invariants, stated as theorems}

\wbq{Check}{Q23}{The consistency-model theorem gives linearizability ``for free.''
What does ``free'' cost? (Hint: the busy-wait and throughput under contention --- the
single writer is also a scalability ceiling.)}

\wbq{Trace}{Q24}{Construct a three-operation claim history and give its unique
linearization. Then show why no second linearization is consistent with real time.}

\wbq{Open \wbstar}{Q25}{\textbf{(OP-3, Cross-runtime soundness as a proof
obligation.)} Specify the minimal differential-test suite that would make I11 a
theorem: what operation sequences, what observable-state assertions, run against both
runtimes?}

% ─── 26–27 · Cross-organ atomicity (paper §12) ────────────────────────────────
\subsection*{Cross-organ atomicity}

\wbq{Check}{Q26}{Give a concrete partial-failure: port claimed, session opened, then
the commitment write fails. What inconsistent state is left, and which agent sees
it?}

\wbq{Open \wbstar}{Q27}{\textbf{(OP-11, Cross-organ atomicity by default.)} Sketch
the interface change that wraps the three-organ ``begin'' operation in one local
transaction. What is the compensating action if you instead chose Sagas, and why is
the single transaction strictly simpler here?}

% ─── 28–30 · Threat model (paper §13) ─────────────────────────────────────────
\subsection*{Threat model and the limits of mediation}

\wbq{Check}{Q28}{The hash chain defends against nobody in scope. Name the
out-of-scope adversary it does defend against, and the layer (substrate or economy)
where that adversary becomes real.}

\wbq{Trace}{Q29}{An agent forges its actor identity to release another agent's lock.
Walk the lock-owner-validity rule and show exactly where I12's absence lets the
forgery through.}

\wbq{Open \wbstar}{Q30}{\textbf{(OP-9, Same-machine adversary.)} What is the minimal
hardening that closes the same-machine adversary without a hardware
trusted-platform-module dependency --- an OS keychain for the signing key, sealed
memory? Where does each stop?}
